\documentclass[10pt,a4paper]{article}
\usepackage[T1]{fontenc}
\usepackage[utf8]{inputenc}
\usepackage{charter}
\usepackage{microtype}
\usepackage[margin=1.7cm,top=1.4cm,bottom=1.4cm]{geometry}
\usepackage{enumitem}
\usepackage{titlesec}
\usepackage{xcolor}
\usepackage{graphicx}
\usepackage[hidelinks]{hyperref}
\definecolor{linkblue}{HTML}{0B57D0}
\definecolor{rule}{HTML}{D9A88A}
\definecolor{accent}{HTML}{C05621}
\hypersetup{colorlinks=true,urlcolor=linkblue,linkcolor=linkblue}

\pagestyle{empty}
\setlength{\parindent}{0pt}

\titleformat{\section}{\normalsize\bfseries\color{accent}}{}{0pt}{}[\vspace{-0.6em}\color{rule}\rule{\textwidth}{0.4pt}]
\titlespacing*{\section}{0pt}{1.0em}{0.5em}

\newcommand{\role}[3]{%
  \textbf{#1}\hfill\textcolor{accent}{\textit{#2}}\hfill{#3}\par\vspace{0.15em}}

\setlist[itemize]{leftmargin=1.1em,itemsep=0.12em,topsep=0.18em,parsep=0pt}

\begin{document}

\begin{minipage}[c]{\dimexpr\textwidth-2.6cm\relax}
{\LARGE\bfseries Giuseppe Baldassarre}\\[0.25em]
{\normalsize Geospatial Analyst \textbar{} Remote Sensing \& Earth Observation}\\[0.15em]
{\small Ph.D. Environmental Engineering (Remote Sensing) \textperiodcentered{} M.Sc. Mediterranean Forestry and Natural Resource Management}\\[0.35em]
{\small Tel. +39 371 361 9920 \textbar{} \href{mailto:gi.baldassarre@gmail.com}{gi.baldassarre@gmail.com} \textbar{} Lisbon, Portugal}\\[0.30em]
{\small\textbf{Portfolio:} \href{https://gbal17.github.io/}{\textbf{gbal17.github.io}} \quad\textperiodcentered\quad \textbf{Full online CV:} \href{https://gbal17.github.io/taleo-cv.html}{\textbf{gbal17.github.io/taleo-cv.html}}}\\[0.15em]
{\footnotesize \href{https://www.linkedin.com/in/giuseppe-baldassarre-671a7623/}{LinkedIn} \textperiodcentered{} \href{https://scholar.google.com/citations?user=M8b10awAAAAJ&hl=en-EN}{Google Scholar} \textperiodcentered{} \href{https://www.researchgate.net/profile/Giuseppe-Baldassarre?ev=prf_overview}{ResearchGate} \textperiodcentered{} \href{https://medium.com/@gi.baldassarre}{Medium} \textperiodcentered{} \href{https://github.com/gbal17}{GitHub}}
\end{minipage}\hfill
\begin{minipage}[c]{2.4cm}
\raggedleft\includegraphics[width=2.3cm]{logo_round.png}
\end{minipage}

\vspace{0.7em}

Remote Sensing and GIS Specialist with over twelve years of professional experience, I design and
operationalize satellite-based monitoring systems for land, water and agriculture. I combine advanced
GIS/RS processing (Python, Google Earth Engine, SEPAL) with decision-support tools and inclusive
governance mechanisms to support FAO, IFAD-funded and UNCCD projects and national partners in
climate-vulnerable countries.

\section{Skills}
\textbf{Languages:} Italian (mother tongue), English (excellent), Spanish (excellent working knowledge), Portuguese and Turkish (intermediate)\\[0.2em]
\textbf{Programming \& data:} Python (GeoPandas, Rasterio, scikit-learn, XGBoost, pandas), Google Earth Engine (JavaScript \& Python API, geemap), MATLAB, Git/GitHub, \LaTeX{}\\[0.2em]
\textbf{Platforms \& tools:} QGIS, PyQGIS, SEPAL, Collect Earth Online, EarthMap, FAO Hand-in-Hand Geospatial Platform, Office\\[0.2em]
\textbf{Domains:} Optical and SAR remote sensing, land cover classification (LCML/ISO 19144-2), change detection, time-series analysis, land degradation (SDG 15.3.1), accuracy assessment, spatial data harmonization and metadata

\section{Professional experience}

\role{Geospatial Analyst}{FAO, Land and Water Division}{2023 -- present}
\begin{itemize}
\item Designed the geospatial indicator framework for the \textbf{IFAD-funded ROOTS project} (Resilience of Organizations for Transformative Smallholder Agriculture, UTF/GAM/047/GAM), The Gambia\,\href{https://openknowledge.fao.org/items/f3175f42-56ff-4b85-892b-528dd124b73a}{(read)}; produced the 2023 Land Cover Map of The Gambia, the first under the LCML/ISO 19144-2 reference system \href{https://openknowledge.fao.org/items/9bb51030-7521-4c72-97d0-5603c5221049}{(read}, \href{https://www.youtube.com/watch?v=5cydX2MixNs}{watch)}, and built the national geospatial monitoring platform now used by government counterparts
\item Leading the GEE-to-PRAIS4 automation pipeline for SDG 15.3.1 / Land Degradation Neutrality reporting across \textbf{21 UNCCD country Parties} (2026 PRAIS4 cycle, COP17 deadline); piloted the full workflow for Mongolia (land cover, LPD, SOC, transition matrices)
\item Designed the district-level INFORM Risk Index for FAO Afghanistan's \href{https://data.apps.fao.org/aware}{AWARE platform} ($\sim$401 districts, 60+ indicators); published 35 composite layers to the CKAN catalogue
\item Produced Syria's 10\,m national land cover map (AlphaEarth Foundations embeddings + Sentinel-2 + Random Forest, GEE) for drought impact assessment \href{https://openknowledge.fao.org/handle/20.500.14283/ce0956en}{(read)}
\item Flood impact assessment on cropland in Pakistan, Oct--Nov 2025 \href{https://openknowledge.fao.org/items/32a5a117-80f8-4038-acd9-040ae5e53a61}{(KP}, \href{https://openknowledge.fao.org/items/b3340185-2f8f-444b-9795-ea3512f774ca}{Punjab}, \href{https://openknowledge.fao.org/items/fb92c557-615a-4c61-a197-c2ebd92b32f0}{Sindh)}
\item Land Cover Atlas of Yemen \href{https://openknowledge.fao.org/items/3ca05f72-f0b0-4ccb-9567-531fa8d84619}{(read)}
\item Development of \href{http://34.78.39.234:8082/}{the Gambia GeoLab} geospatial platform, presented at the \href{https://earthoutreachonair.withgoogle.com/events/geoforgood25-singapore}{Geo for Good Summit} in Singapore, 2025 \href{https://www.fao.org/geospatial/news/news-detail/fao-presents-the-gambia-geolab-at-geo-for-good-summit-in-singapore/en}{(FAO news}, \href{https://medium.com/@gi.baldassarre/turning-research-and-data-into-actionable-insights-for-agriculture-and-natural-resource-management-b2d035aa96a8}{read)}
\item Design and delivery of trainings on ``Geospatial Tools for Mangrove Monitoring in The Gambia'' \href{https://drive.google.com/file/d/1bKH9oAmR6DfebORUtqPH44YM-akl8gxG/view}{(visit)} and ``Using SEPAL, Google Earth Engine and Collect Earth Online to produce a Land Cover Map of Gambia (2023)'' \href{https://drive.google.com/file/d/1gtUpSWBR5gW9FGgZv1lKpaMOALRHdjy3/view}{(visit)}
\end{itemize}

\role{Post-doctoral researcher}{University of Lisbon, Portugal}{2020 -- 2023}
\begin{itemize}
\item Estimating shrub live fuel moisture content for mainland Portugal \href{https://www.mdpi.com/2571-6255/8/5/178}{(read)}
\item Study of traditional extensive pastoralism relying on fires in Northern Portugal with GIS/RS analysis
\item Mapping fire regimes over Portugal \href{https://www.researchgate.net/publication/366124207_Regimes_de_fogo_em_Portugal_Continental}{(read1}, \href{https://www.researchgate.net/publication/360227465_Cartografia_de_Regimes_de_Fogo_a_Escala_da_Freguesia_1980-2017}{read2)} and Turkey \href{https://www.researchgate.net/publication/389813750_Describing_fire_regimes_over_Turkey_using_MODIS_fire_observations}{(read)}
\item Automatic monitoring of the national fuel break network of Portugal with Sentinel-2 in GEE \href{https://link.springer.com/chapter/10.1007/978-3-030-96466-5_5}{(read)}
\end{itemize}

\role{EU Aid Volunteer}{Sustainable Rural Development, Viet Nam}{2021 -- 2022}
\begin{itemize}
\item SAR-based above-ground biomass and carbon stock estimation for mangrove restoration (VM069), with a \href{https://baldassarre.users.earthengine.app/view/srd}{GEE application} to monitor mangrove extent in Southern Viet Nam
\item Developing an \href{https://ee-gibaldassarre.projects.earthengine.app/view/tccvn}{App} in GEE to monitor tree canopy cover in Viet Nam (2000--2021)
\item Mapping smallholder coffee production systems driving land use change in NW Viet Nam \href{https://www.youtube.com/watch?v=i9aKtECsYz8}{(watch)}
\end{itemize}

\role{Post-doctoral researcher}{Istanbul Technical University, Turkey}{2012 -- 2018}
\begin{itemize}
\item Biomass burning emission estimation with geostationary observations \href{https://acp.copernicus.org/articles/15/8539/2015/acp-15-8539-2015.html}{(read)}
\item Description of forest fires over Turkey using the most advanced satellite observations \href{https://fires.itu.edu.tr/bbeurasia/index.html}{(visit)}
\item First version of the national Forest Fire Monitoring System web-GIS in Turkey (UYANBIL) \href{https://fires.itu.edu.tr/}{(visit)}
\end{itemize}

\role{Environmental consultant (volunteer)}{Aprode Peru (NGO), Peru}{2011}
\begin{itemize}
\item Field campaigns to outline the solid waste management cycle in rural villages in Northern Peru
\item Implemented a project for the social inclusion of informal recyclers \href{https://www.youtube.com/watch?v=av93xA8a-ls}{(watch)}
\end{itemize}

\role{Visiting scholar}{University of Wisconsin--Madison, USA}{2008, 2009 -- 2010}
\begin{itemize}
\item Development and assessment of advanced satellite fire detection products using statistical tools
\item Testing product operational applicability over Southern Italy
\item Publications: \href{https://www.researchgate.net/publication/234055644_Automatic_RST-based_system_for_a_rapid_detection_of_fires}{``Automatic RST-based system for a rapid detection of fires''}; \href{https://www.researchgate.net/publication/306536266_Assessment_of_the_Robust_Satellite_Technique_RST_in_real_time_detection_of_summer_fires}{``Assessment of the Robust Satellite Technique (RST) in real-time detection of summer fires''}
\end{itemize}

\section{Education}

\role{M.Sc., Mediterranean Forestry and Natural Resource Management}{\href{https://www.medfor.eu/}{Erasmus Mundus MEDfOR}}{2018 -- 2020}
\begin{itemize}
\item 1st year: University of Lisbon (Portugal), Valladolid (Spain) \& Padova (Italy) --- natural resource management with decision support systems, forest governance, agroforestry, fire ecology \& management, GIS, research \& project development methodology
\item 2nd year: University of Porto (Portugal) \& Lleida (Spain) --- governance of Mediterranean forests, environmental economics, research methods
\item Thesis: \href{https://www.researchgate.net/publication/389813750_Describing_fire_regimes_over_Turkey_using_MODIS_fire_observations}{``Describing fire regimes over Turkey using MODIS fire observations''}. Graduated with excellent standing (8.91/10)
\end{itemize}

\role{Ph.D., Environmental Engineering (Remote Sensing)}{U. of Wisconsin--Madison \& U. of Basilicata}{2007 -- 2011}
\begin{itemize}
\item Advanced tools for satellite-based environmental monitoring (biomass combustion emissions)
\item Thesis: ``Performance of the most advanced satellite techniques for real-time monitoring of forest fires over Italy''
\item Publication: \href{https://www.researchgate.net/publication/234245164_A_Total_Validation_Approach_for_assessing_the_RST_technique_in_forest_fire_detection_and_monitoring}{``A Total Validation Approach for assessing the RST technique in forest fire detection and monitoring''}
\end{itemize}

\role{M.Sc., Environmental Engineering (110/110 \textit{summa cum laude})}{University of Basilicata, Italy}{2004 -- 2007}
\begin{itemize}
\item Thesis: ``Robust satellite techniques for fire monitoring using geostationary satellites (MSG-SEVIRI)''
\item Waste-water and waste management; atmospheric physics; water, air and soil pollution dynamics
\item Publication: \href{https://www.researchgate.net/publication/234055648_Assessment_of_a_Robust_Satellite_Technique_for_forest_fire_detection_and_monitoring_by_using_a_Total_Validation_Approach_TVA}{``Assessment of a Robust Satellite Technique for forest fire detection and monitoring using a Total Validation Approach (TVA)''}
\end{itemize}

\role{B.Sc., Environmental Engineering (110/110 \textit{summa cum laude})}{University of Basilicata, Italy}{2001 -- 2004}
\begin{itemize}
\item Physics, chemistry, geology, statistics, economics, applied environmental engineering
\end{itemize}

\vspace{0.6em}
\small Supporting documents: \href{https://drive.google.com/file/d/1TurHo7zIqE6_efP7Ho0fWwewve0lyJq2/view?usp=sharing}{certificates}, \href{https://drive.google.com/file/d/1QmVIYiWMHVG7rS61SMaESllyqTxkrcRU/view}{grades}, \href{https://drive.google.com/file/d/1D3z3y6sKUg0MaK5rGoZvrykGRFlAcdOk/view?usp=share_link}{extra-courses}, \href{https://drive.google.com/file/d/1KgUw_HTc51RGX6VVkrUTtVh2zfrgpXEz/view?usp=share_link}{international cooperation}, \href{https://drive.google.com/file/d/1dj5vOcAGKw-qrIa63V-UigJ2MN-MjlRj/view}{references}.

\end{document}
